% Dynamic simulation chapter skeleton. Requires \TemplatePrompt.
\section{动态系统与仿真分析}

\subsection{系统边界、状态与时间尺度}
\TemplatePrompt{定义状态量、控制量、外生输入、时间步长、初值、边界条件和守恒关系。}

\subsection{可运行基线}
\TemplatePrompt{采用无干预、稳态、解析近似或简化动力学作为同输出基线，并保持相同初值和时域。}

\subsection{状态方程与参数来源}
\TemplatePrompt{给出连续或离散状态方程、耦合关系、参数单位、标定数据和近似假设。}

\subsection{数值方法与误差控制}
\TemplatePrompt{记录积分器、步长/容差、事件处理、随机种子、收敛检查和计算成本。}

\subsection{轨迹、情景与机制解释}
\TemplatePrompt{比较主模型与基线的轨迹、阈值、峰值和累积量；解释差异来自哪个动力学机制。}

\subsection{稳定性、扰动、稳健性与结论边界}
\TemplatePrompt{执行步长收敛、参数扰动、极端初值和外生冲击测试；报告结论的稳健性与改变条件。}
